\PlanPage{5}{Section 3, Methodology: 6 pages, pp.~5--10. Section 3.1: 1 page, p.~5.}
\section{Methodology}\label{sec:framework}
% First subsection of Methodology; included by framework.tex.
\subsection{Problem Formulation}\label{sec:problem}
\subsubsection{Target Flows and Observable Evidence}
\Budget{0.50 page}
\DraftText{Define a target flow, its binary label, protocol messages, endpoint
roles, observation cutoff, and the allowed related-flow history. State which
protocol fields are observable. Specify what evidence is available when a
decision is issued and how packet/flow identifiers are preserved. Relate the
task to heterogeneous traffic without assuming visibility into encrypted payloads.}
\DraftBlock{3.8cm}{Notation and one worked target-flow example:\\
target flow, source messages, endpoint roles, time window, context budget, output label.}
\subsubsection{Decision Paths and Processing Budget}
\Budget{0.50 page}
\DraftText{Define the complete input stream, the fast detector's alert set, and
the independently selected discovery candidates. Both selections originate from
the full stream; the discovery set is not restricted to alerts. Define the
final prediction unit, candidate overlap, semantic invocation count, and context
budget. State the detection and processing objectives without assigning
unmeasured numerical gains.}
\DraftBlock{3.2cm}{Task specification and processing constraints.\\
Freeze actual thresholds, observation rules, and evaluation granularity before experiments.}


\PlanPage{6}{Section 3.2: 1 page, p.~6; part 2 of the 6-page method allocation.}
\subsection{System Overview}
\Budget{1 page}
\DraftText{Introduce the two semantic layers and two parallel decision paths
through their inputs and outputs. Walk one target flow from packet observation
to its final prediction. Associate each module with the challenge it resolves.}
\DraftFigure{7.0cm}{Overview of the proposed framework}{Figure 1 layout:\\[6pt]
Full traffic stream $\longrightarrow$ fast scoring $\longrightarrow$ alert verification\\[6pt]
Full traffic stream $\longrightarrow$ independent routing $\longrightarrow$ cross-flow discovery\\[6pt]
Shared evidence construction: message-level operational semantics + cross-flow behavioral semantics\\[6pt]
Jev typed decisions $\longrightarrow$ integration $\longrightarrow$ target-flow predictions.\\[6pt]
Show packet/flow provenance and candidate overlap handling.}
\DraftText{Use the architecture to establish the responsibilities of the
modules in Sections 3.3--3.7. Provide a concise execution description or
pseudocode for the actual implementation.}

\PlanPage{7}{Section 3.3: 1 page, p.~7; part 3 of the 6-page method allocation.}
\subsection{Message-Level Operational Semantics}\label{sec:message-semantics}
\Budget{1 page}
\DraftText{Define the common operation, resource, response-outcome, and
endpoint-role slots. Describe how implemented protocol parsers populate them,
how unavailable fields are represented, and how source packet identifiers are
retained. Use the same operational intent in multiple supported protocols to
explain alignment.}
\DraftTable{5.0cm}{Examples of protocol-to-operation mapping}{Protocol / source field / operation / resource / response outcome / endpoint role / provenance.\\[6pt]
Populate with implemented mappings and trace-backed examples.}
\DraftText{Explain the semantic invariants and protocol-specific mapping
rules. Summarize the mapping schema here; put the extended mappings in
Appendix~\ref{app:reproducibility}.}
\DraftBlock{3.2cm}{Mapping rules, missing-value handling, and an end-to-end message example.}

\PlanPage{8}{Section 3.4: 1 page, p.~8; part 4 of the method allocation.}
\subsection{Cross-Flow Behavioral Semantics}\label{sec:cross-flow}
\Budget{1 page}
\DraftText{Start from the target flow. Define endpoint and temporal
association, event order, and the selected communication history. Explain the
deterministic counts, rates, and response differences. Describe template folding
and the preservation of meaningful changes, followed by the exact context
truncation rule and provenance mapping.}
\DraftFigure{6.0cm}{Construction of bounded cross-flow evidence}{Related communications in time order $\longrightarrow$ behavioral facts\\[6pt]
$\longrightarrow$ repeated-message folding $\longrightarrow$ bounded Jev state.\\[6pt]
Annotate retained evidence, changed outcomes, source IDs, and token budget.}
\DraftBlock{3.8cm}{Algorithm and illustrative context before/after folding.\\
State the association and ordering rules used by the implemented system.}

\PlanPage{9}{Section 3.5: 1 page, p.~9; two half-page path descriptions.}
\subsection{Dual-Path Candidate Selection}\label{sec:paths}
\subsubsection{Alert Verification}
\Budget{0.50 page}
\DraftText{Describe full-stream fast scoring, the resulting alert set, and
the evidence used by Jev to reassess each selected target. Specify actual
selection and decision thresholds. Explain how verification changes the final
target-flow prediction and how its precision effect will be measured.}
\DraftBlock{3.8cm}{Verification path: candidate rule, typed decision, retained/revised alert.}
\subsubsection{Cross-Flow Discovery}
\Budget{0.50 page}
\DraftText{Define the association cues and routing rule applied independently
to the full stream. Explain why selected target flows need related-flow evidence
and how discovery produces additional target-flow decisions. Describe duplicate
handling when a target is selected by both paths.}
\DraftBlock{3.8cm}{Discovery path: full-stream candidates, evidence construction, additional decisions.}

\PlanPage{10}{Sections 3.6--3.7: 1 page, p.~10; part 6 of the method allocation.}
\subsection{Typed Decisions and Output Integration}\label{sec:typed}
\Budget{0.50 page}
\DraftText{Specify the serialized state, typed question, answer type, and
probability returned by the existing Jev model~\cite{R11,R02}. Explain the
implemented calibration and integration rule for the fast detector and both
semantic paths, including overlapping candidates and failed requests. Freeze
the rule from the actual implementation; do not insert an earlier unconfirmed
fusion formula.}
\DraftBlock{3.7cm}{Typed interface and final output rule.\\
Model identifier/version: xxx. Thresholds and overlap rule: xxx.}
\subsection{Implementation and Reproducibility}
\Budget{0.50 page}
\DraftText{Give parser and dependency versions, streaming/window management,
hardware, API/client configuration, concurrency, timeout/retry handling, and
component timing boundaries. Identify the code revision and available artifacts.
Record applicable AI assistance in preparing text, figures, or study design.
The extended schema and run manifest are in Appendix~\ref{app:reproducibility}.}
\DraftBlock{3.3cm}{Implementation details and reproducibility pointers.\\
Software versions, code revision, model snapshot, and configuration IDs: xxx.}
